% BHU PYQ -- Manually transcribed & error-corrected
% Paper: BCH-122 : Financial Accounting - II
% Exam: B.Com. (Hons.) Semester II Examination, 2014-15 (Special)

\documentclass[11pt,a4paper]{article}
\usepackage[a4paper,top=2.5cm,bottom=2.5cm,left=2.8cm,right=2.8cm,headheight=14pt]{geometry}
\usepackage{amsmath,amssymb}
\usepackage[T1]{fontenc}
\usepackage[utf8]{inputenc}
\usepackage{lmodern,microtype,fancyhdr,enumitem,hyperref,lastpage,parskip}

\hypersetup{
    colorlinks=true,
    urlcolor=black,
    linkcolor=black,
    pdftitle={BCH-122: Financial Accounting - II -- BHU B.Com. Sem II 2014-15}
}

\pagestyle{fancy}
\fancyhf{}
\renewcommand{\headrulewidth}{0.4pt}
\renewcommand{\footrulewidth}{0.4pt}
\fancyhead[L]{\small\textit{Banaras Hindu University}}
\fancyhead[R]{\small\textit{BCH-122: Financial Accounting - II}}
\fancyfoot[C]{\small Page \thepage\ of \pageref{LastPage}}

\newlist{parts}{enumerate}{2}
\setlist[parts,1]{label=(\alph*),leftmargin=*,itemsep=5pt,topsep=3pt}
\setlist[parts,2]{label=(\roman*),leftmargin=*,itemsep=3pt,topsep=2pt}
\newcommand{\pts}[1]{\hfill\textit{[#1]}}

\begin{document}

\begin{center}
  {\large\bfseries BANARAS HINDU UNIVERSITY}\\[2pt]
  {\normalsize B.Com. (Hons.) --- Semester II (Special) Examination, 2014-15}\\[6pt]
  \rule{0.8\linewidth}{0.5pt}\\[6pt]
  {\large\bfseries Paper No. BCH-122: Financial Accounting - II}\\[4pt]
  {\normalsize Time: 3 Hours\hfill Maximum Marks: 70}
\end{center}

\rule{\linewidth}{0.4pt}

\smallskip
\noindent\textit{Instructions: Answer all questions. Marks are indicated against each question.}

\bigskip

\noindent\textbf{Question 1.}
\begin{parts}
  \item How are Bill of Exchange and Promissory Note defined in the Negotiable Instruments Act? Describe the distinguishing features of the two.\pts{6}
  \item On 30th December 2014, the account of $X$ according to the Bank column of the Cash Book was overdrawn to the extent of Rs. 4,062. On the same date, the Bank statement showed a balance of Rs. 1,400 in favor of $X$. An examination of the Cash Book and Bank statement revealed the following:
  \begin{parts}
    \item A cheque for Rs. 1,140 deposited on 29th September 2014 was credited by the bank on 3rd October 2014.
    \item A payment by cheque for Rs. 160 has been entered twice in the Cash Book.
    \item On 29th September 2014, the bank credited an amount of Rs. 1,740 received from a customer of $X$, but the advice was not received by $X$ until 1st October 2014.
    \item Bank charges amounting to Rs. 58 had not been entered in the Cash Book.
    \item On 6th September 2014, the Bank credited Rs. 2,000 to $X$ in error.
    \item A Bill of Exchange for Rs. 1,000 was discounted by $X$ with his Bank. This bill was dishonoured on 28th September 2014, but no entry was made in the books of $X$.
    \item Cheques issued upto 30th September 2014 but not presented for payment upto that date totalled Rs. 3,760.
  \end{parts}
  You are required to:
  \begin{parts}
    \item Show the appropriate rectifications required in the Cash Book of $X$ to arrive at the correct balance on 30th September 2014.
    \item Prepare a Bank Reconciliation Statement as on that date.\pts{8}
  \end{parts}
\end{parts}

\centerline{\textbf{OR}}

\begin{parts}
  \item What is a Bank Reconciliation Statement? When and why is it prepared?\pts{6}
  \item On January 1st 2014, $M$ accepts two bills for Rs. 5,000 and Rs. 6,000 respectively, for the amount payable to $P$. The first bill is payable after two months and the second bill is payable after three months.
  
  On February 4, 2014, $P$ gets the second bill discounted from the bank at a discount of Rs. 400. On February 15, 2014, the first bill of Rs. 5,000 is endorsed in favor of $R$ in full and final settlement of Rs. 5,500 due to him.
  
  On due dates, the first bill is duly honoured but the second bill is dishonoured and the bank pays Rs. 50 as noting charges. On April 10, 2014, $M$ pays Rs. 600 for interest payable on renewal of the bill and accepts a new bill payable after sixty days. On May 1, 2014, $M$ is declared insolvent and a final dividend of 20\% is received from his estate.
  
  Pass journal entries in the books of $P$ and $M$ and show the account of $P$ in the books of $M$.\pts{8}
\end{parts}

\medskip\rule{\linewidth}{0.2pt}

\newpage

\noindent\textbf{Question 2.}
\begin{parts}
  \item Explain the distinguishing features of the types of shares which can be issued by a company.\pts{6}
  \item A company issued for public subscription 10,000 shares of Rs. 10 each at Rs. 11 per share payable as follows:
  \begin{itemize}[label=-,leftmargin=1.5cm]
    \item Rs. 3 on application
    \item Rs. 4 on allotment (including premium)
    \item Rs. 4 on call
  \end{itemize}
  Applications were received for 12,000 shares and the directors made pro-rata allotment. 
  
  $A$, an applicant for 120 shares, could not pay the allotment and call money. $B$, a holder of 200 shares, failed to pay the call. All those shares were later on forfeited. Out of the forfeited shares, 150 shares (the whole of $A$'s shares being included) were reissued at Rs. 9 per share credited as fully paid.
  
  Prepare Cash Book and pass the journal entries in the books of the company. Also show the Balance Sheet of the company.\pts{8}
\end{parts}

\centerline{\textbf{OR}}

\begin{parts}
  \item Describe the provisions of Section 80 of the Companies Act as regards redemption of preference shares.\pts{6}
  \item The Balance Sheet of a company on January 1, 2014 was as follows:
  \begin{center}
    \small
    \begin{tabular}{|l|r||l|r|}
      \hline
      \textbf{Liabilities} & \textbf{Amount (Rs.)} & \textbf{Assets} & \textbf{Amount (Rs.)} \\
      \hline
      Equity share capital (Rs. 10 each) & 2,00,000 & Sundry Assets & 5,00,000 \\
      Preference share capital (Rs. 100 each) & 1,00,000 & & \\
      Securities Premium & 8,000 & & \\
      General Reserve & 40,000 & & \\
      Profit \& Loss A/c & 50,000 & & \\
      Other liabilities & 1,02,000 & & \\
      \hline
      \textbf{Total} & \textbf{5,00,000} & \textbf{Total} & \textbf{5,00,000} \\
      \hline
    \end{tabular}
  \end{center}
  Preference shares were due for immediate redemption at 10\% premium. To finance the redemption, investments of the book value of Rs. 40,000 were sold at a profit of 25\% and equity shares were issued at a premium of 20\%. It was decided to utilise the free reserves to the minimum extent possible.
  
  On 31.03.2014, it was decided to issue bonus shares in the ratio of 2 for 5.
  
  Pass journal entries and show the Balance Sheet after the bonus issue, assuming that the holder of 100 preference shares is not traceable.\pts{8}
\end{parts}

\medskip\rule{\linewidth}{0.2pt}

\newpage

\noindent\textbf{Question 3.}
Distinguish between a share and a debenture. Describe the methods of and accounting for redemption of debentures.\pts{14}

\centerline{\textbf{OR}}

\begin{parts}
  \item Distinguish between Capital Redemption Reserve and Debenture Redemption Reserve.\pts{6}
  \item The Balance Sheet of $Z$ Ltd. disclosed the following information on 1.1.2013:
  \begin{itemize}[label=-,leftmargin=1.5cm]
    \item 13\% Debenture Account: Rs. 7,00,000
    \item Debenture Redemption Fund A/c: Rs. 5,00,000
    \item 13\% Debenture Redemption Fund Investment A/c: Rs. 5,00,000
  \end{itemize}
  The Annual contribution to the Debenture Redemption Fund was Rs. 70,000 for the year 2013 and 2014. The debentures were redeemable on 31st December 2014. On 31st December 2014 the investments were sold for Rs. 7,00,000 and the debentures were redeemed.
  
  Prepare Debenture Account, Debenture Redemption Fund Account and Debenture Redemption Investment Account for the years 2013 and 2014.\pts{8}
\end{parts}

\medskip\rule{\linewidth}{0.2pt}

\noindent\textbf{Question 4.}
How is purchase consideration ascertained in case of acquisition of business by a company? What accounting is done in the books of the company on such purchase? Illustrate.\pts{14}

\centerline{\textbf{OR}}

\begin{parts}
  \item Discuss the basis on which profit prior to and after incorporation of a company is separated.\pts{6}
  \item Triveni Limited was incorporated on May 1, 2014 to purchase the business of Triveni Sahay as from January 1, 2014 and obtained its certificate for commencement of business on June 1, 2014.
  
  The accounts of the company for the period ended December 31, 2014 disclosed the following facts:
  \begin{parts}
    \item The Turnover for the whole year amounted to Rs. 2,40,000 of which Rs. 40,000 related to the period from January 1, 2014 to April 30, 2014.
    \item The Trading Account disclosed a Gross Profit of Rs. 96,000.
    \item The following items appeared in the Profit and Loss Account:
    \begin{center}
      \small
      \begin{tabular}{|l|r||l|r|}
        \hline
        \textbf{Particulars} & \textbf{Amount (Rs.)} & \textbf{Particulars} & \textbf{Amount (Rs.)} \\
        \hline
        Director's Fee & 1,500 & General Expenses & 1,800 \\
        Auditor's Fee & 750 & Commission on Sales & 3,600 \\
        Rent \& Taxes & 4,800 & Printing \& Stationery & 2,400 \\
        Bad Debts (Rs. 700 post-incorporation) & 2,000 & Advertisement & 4,200 \\
        Salaries & 12,000 & Travelling Exp. & 8,400 \\
        Depreciation & 3,600 & Interest to vendors & 4,000 \\
        Debenture Interest & 6,000 & (6\% p.a. on Rs. 1,00,000) & \\
        Preliminary Expenses & 2,400 & & \\
        \hline
      \end{tabular}
    \end{center}
  \end{parts}
  From the above, find out the profit available for dividend and the profit prior to incorporation.\pts{8}
\end{parts}

\medskip\rule{\linewidth}{0.2pt}

\newpage

\noindent\textbf{Question 5.}
\begin{parts}
  \item Distinguish between charges against profit and appropriation of profit.\pts{6}
  \item Following are the balances in the Balance Sheet of a company of limited liability as on June 30, 2014:
  \begin{parts}
    \item Authorised capital: 25,000 shares of Rs. 100 each.
    \item Issued capital: 20,000 shares, Rs. 90 per share called and paid up.
    \item Reserve Fund: Rs. 9,00,000.
    \item Profit \& Loss Account: Rs. 4,00,000.
  \end{parts}
  The following resolutions were passed in the Annual General Meeting of the company held on August 18, 2014:
  \begin{parts}
    \item To pay Rs. 15 each dividend for each share.
    \item To declare a capital Bonus @ Rs. 50 per share to be used in making partly paid shares as fully paid and the balance to be used in issuing fully paid bonus shares at a premium of 60\%.
    \item Rs. 50,000 to be paid as bonus to the employees of the company.
  \end{parts}
  These resolutions were carried out. Pass the necessary journal entries in the books of the company.\pts{8}
\end{parts}

\vfill
\rule{\linewidth}{0.4pt}
\begin{center}\small
\href{https://bhu-exam.vercel.app/}{Click here to visit our website}\\[4pt]
\href{https://me-aryan.vercel.app/}{Click here to visit our contributor}\\[4pt]
\href{https://quashan-me.vercel.app/}{Click here to visit our contributor}
\end{center}

\end{document}
